\documentclass[11pt]{article}
\usepackage[letterpaper,margin=1in]{geometry}
\usepackage[T1]{fontenc}
\usepackage{lmodern,microtype,amsmath,amssymb,amsthm}
\usepackage{xcolor}
\usepackage[colorlinks=true,linkcolor=black,citecolor=blue!45!black,urlcolor=blue!45!black]{hyperref}
\usepackage{xurl}
\newtheorem{theorem}{Theorem}[section]
\newtheorem{lemma}[theorem]{Lemma}
\newtheorem{corollary}[theorem]{Corollary}
\theoremstyle{definition}
\newtheorem{definition}[theorem]{Definition}
\theoremstyle{remark}
\newtheorem{remark}[theorem]{Remark}
\newcommand{\vol}{\mathrm{vol}}
\newcommand{\F}{\mathbb{F}}
\newcommand{\AEST}{\mathsf{AEST}}
\newcommand{\TSUM}{\mathsf{3SUM}}
\newcommand{\Inc}{\mathrm{Inc}}
\newcommand{\id}{\mathrm{id}}
\newcommand{\ind}[1]{[\,#1\,]}
\hypersetup{pdftitle={Superlinear circuit lower bounds for All-Edges Sparse Triangle, 3SUM, and shortest paths},
  pdfauthor={Charlie Cruz},
  pdfkeywords={circuit lower bounds, All-Edges Sparse Triangle, APSP, 3SUM, triangle detection, element distinctness, tensor networks, unambiguous verification, superlinear lower bounds}}
\title{Superlinear Circuit Lower Bounds for All-Edges Sparse Triangle,\\
3SUM, and Shortest Paths, Assuming a Tensor-Network Volume Bound}
\author{Charlie Cruz}
\date{November 25, 2026}
\begin{document}
\maketitle

\begin{abstract}
We prove superlinear Boolean circuit lower bounds for three central
problems in fine-grained complexity, assuming the tensor-network volume
lower bound for list distinctness stated as Theorem~2.2 of \emph{The matrix
multiplication exponent is nine fourths}. Our main result is that
All-Edges Sparse Triangle, where one must decide for every edge of an
$m$-edge graph whether it lies in a triangle, requires
$\Omega(m(\log m)^{1+\sigma})$ gates for a constant $\sigma>0$. This holds
even for sorted edge lists of 2-degenerate graphs on vertex set
$\{0,\dots,m-1\}$, and the same bound applies to counting triangles and to
deciding whether one exists. As corollaries, answering $m$ distance
queries in a sparse unweighted DAG, and hence APSP on such graphs, requires
the same number of gates. We also show that 3SUM on $n$ integers requires
$\Omega(n(\log n)^{1+\sigma})$ gates. Each bound is superlinear in the input
length in bits. The triangle lower bound comes from two simple gadgets: a
graph whose triangles perform parallel table lookups, and a sorting
certificate that three lookups can check and that is unique exactly when
the input list is distinct.
\end{abstract}

\section{Introduction}

Three problems sit at the heart of fine-grained complexity: 3SUM,
All-Pairs Shortest Paths (APSP), and triangle detection. Each is believed
to require essentially its textbook running time, and a large web of
reductions ties their complexities to those of many other problems. In
particular, \emph{All-Edges Sparse Triangle} ($\AEST$), where one is given
a graph with $m$ edges and must determine for every edge whether it lies in
a triangle, is known to require $m^{4/3-o(1)}$ time on a word RAM if either
the 3SUM hypothesis or the APSP hypothesis holds \cite{VWX}.

These lower bounds are conditional, and they concern a specific
computational model. Unconditionally, much less is known. Against general
Boolean circuits, until very recently the best lower bound for any
explicit function was of the form $cn$ for a small constant $c$, and no
superlinear bound was known for any problem in $\mathsf P$ (see the
discussion in \cite{R06}).

The recent paper \emph{The matrix multiplication exponent is nine fourths}
\cite{OAI} states, as its Theorem~2.2, a superlinear lower bound on the
volume of tensor networks that test whether a list of words is distinct.
Result~06 of this collection \cite{R06} observed that this theorem implies
that element distinctness on $N$ input bits requires
$\Omega(N(\log N)^{\sigma})$ gates. In this paper, we show that the same
theorem gives superlinear circuit lower bounds for all three of the
problems above.

\subsection{Our results}

Throughout, circuits have fan-in at most two, use arbitrary binary gates,
and have unrestricted depth and fan-out; size counts gates. All of our
results assume Theorem~\ref{thm:import} below, which is the single-list
case of \cite[Theorem~2.2]{OAI} over $\F_2$, and $\sigma>0$ always denotes
its exponent. No 3SUM, APSP, or ETH hypothesis is used.

An $\AEST$ instance of size $m$ is a list of $m$ edges $(a,b)$ with
$0\le a<b<m$, each endpoint written in $L_m=\lceil\log m\rceil$ bits. The
input has $B_m=2mL_m=\Theta(m\log m)$ bits, and the output is one bit per
edge, in input order. Our main result is the following.

\begin{theorem}\label{thm:main}
For every sufficiently large $m$, every circuit for $\AEST$ on $m$ edges
has $\Omega(m(\log m)^{1+\sigma})$ gates, equivalently
$\Omega(B_m(\log B_m)^{\sigma})$ gates. This holds even if the circuit need
only be correct on graphs that
\begin{enumerate}
\item are simple, with exactly $m$ vertices $\{0,\dots,m-1\}$ and $m$ edges;
\item are given as lexicographically sorted edge lists; and
\item are 2-degenerate.
\end{enumerate}
The circuit may share gates arbitrarily among its $m$ outputs.
\end{theorem}

The hard instances are thus about as structured as one could hope for.
The same approach handles the natural variants of triangle detection.

\begin{corollary}\label{cor:variants}
Each of the following requires $\Omega(m(\log m)^{1+\sigma})$ gates on
sorted edge lists of simple 2-degenerate graphs with $m$ edges:
\begin{enumerate}
\item counting the triangles, on $m$ vertices;
\item returning, for every edge, a triangle containing it or a marker that
none exists, on $m$ vertices;
\item deciding whether a triangle exists, on $m+1$ vertices; and
\item returning some triangle or a marker that none exists, on $m+1$
vertices.
\end{enumerate}
\end{corollary}

Since $\AEST$ reduces to distance queries, Theorem~\ref{thm:main} also
gives a lower bound for shortest paths.

\begin{corollary}\label{cor:apsp}
For every sufficiently large $m$, every circuit that, given an unweighted
directed acyclic graph with $O(m)$ vertices and $4m$ edges and $m$ query
pairs $(s_k,t_k)$, outputs the distances $d(s_k,t_k)$ has
$\Omega(m(\log m)^{1+\sigma})$ gates. This holds even if only the bits
$\ind{d(s_k,t_k)=2}$ are required. In particular, every circuit for APSP
on such graphs has $\Omega(m(\log m)^{1+\sigma})$ gates.
\end{corollary}

Finally, we show that 3SUM requires superlinear circuits as well. Let
$\TSUM_n$ take $n$ integers of magnitude at most $n^4$ and decide whether
three of them, at distinct positions, sum to zero.

\begin{theorem}\label{thm:3sum-intro}
For every sufficiently large $n$, every circuit for $\TSUM_n$ has
$\Omega(n(\log n)^{1+\sigma})$ gates, even if it need only be correct on
lists of pairwise distinct integers.
\end{theorem}

In every case the input has $B$ bits for $B$ equal to the stated parameter
times a logarithm, so each bound reads $\Omega(B(\log B)^{\sigma})$: it is
superlinear in the input length, by a factor of $(\log B)^\sigma$.

\subsection{Technique overview}

\paragraph{Triangles as lookups.} Our starting point is that triangles in
a sparse graph can implement table lookup. Given an address array
$j\in[q]^q$, a claimed-value array $x\in[q]^q$, and a table $y\in[q]^q$,
consider the predicate $R(j,x,y)=\bigwedge_i\ind{x_i=y_{j_i}}$. We build a
graph on three classes $I$, $A$, $B$ of size $q$: vertex $I_i$ is joined to
$A_{j_i}$ and $B_{x_i}$, and $A_a$ is joined to $B_{y_a}$. Since $I_i$ has
only these two neighbors, the edge $I_iA_{j_i}$ lies in a triangle exactly
when $x_i=y_{j_i}$. So $q$ outputs of an $\AEST$ circuit compute $R$. The
graph has $3q$ vertices and $3q$ edges, its sorted edge order does not
depend on the data, and writing it down costs no gates at all.

\paragraph{A unique sorting certificate.} The predicate $R$ is easy, so we
need a way to build a hard function out of a few lookups. The key idea is
to use nondeterminism of a kind that tensor networks can absorb. Given a
list $X_0,\dots,X_{q-1}$, a certificate consists of a rank array $j$, a
claimed inverse $z$, and a claimed sorted list $Y$. One lookup checks
$z_{j_i}=i$, which forces $j$ to be a permutation; two more check
$Y_{j_i}=X_i$, which forces $Y$ to be a rearrangement of $X$; and a
comparator chain of linear size checks that $Y$ is strictly increasing.
If $X$ is distinct then exactly one certificate passes, and otherwise none
does.

\paragraph{Summing out the certificate.} A circuit is a tensor network of
linear volume \cite{OAI,R06}. Summing over the certificate wires counts
the accepting certificates, and because that count is always $0$ or $1$,
the result is exactly the distinctness indicator, even over $\F_2$. This
summation is a tensor contraction rather than an enumeration, so it costs
only the volume of the certificate wires. Comparing with the volume lower
bound of \cite{OAI} proves Theorem~\ref{thm:main}. For triangle detection
(Corollary~\ref{cor:variants}), where only one output bit is available, we
recover $R$ by a second contraction over selector bits that switch
individual triangles on and off.

\paragraph{Shortest paths and 3SUM.} For Corollary~\ref{cor:apsp}, an edge
$\{u,v\}$ lies in a triangle exactly when $u$ and $v$ have a common
neighbor, which is a distance query in a three-layer copy of the graph.
3SUM is different. The standard reductions go \emph{from} 3SUM and APSP
\emph{to} $\AEST$ \cite{VWX}, which is the wrong direction for
transferring a lower bound, and Theorem~\ref{thm:main} does not seem to
imply Theorem~\ref{thm:3sum-intro}. Instead, we reduce distinctness to
3SUM directly. The difficulty is to detect equal values only at
\emph{different} positions, and we handle it by writing each position into
the low-order bits of the corresponding integers.

\subsection{Discussion}

\paragraph{Comparison with direct reductions.} A direct reduction from
distinctness to triangle detection must compare words at different
positions, which takes a gadget on all pairs $i<j$. Paying for that gadget
forces one to batch $\Theta(q)$ lists, using the many-list form of
\cite[Theorem~2.2]{OAI}, and to name vertices from a universe of size
$\Theta(m^{3/2})$. The certificate approach avoids both. It uses a single
list, vertex names in $[m]$, and graphs of degeneracy $2$. It is an
instance of the observation in \cite[Section~7]{R06} that unambiguous
verifiers convert into tensor networks.

\paragraph{What the bounds do not say.} Our bounds are Boolean circuit
lower bounds. They are not word-RAM bounds, and they are far from the
conditional $m^{4/3-o(1)}$ and $n^{2-o(1)}$ time bounds. Indeed, on our
hard triangle instances, the designated answers can be computed with $q$
indexed table accesses and comparisons, which is $O(m)$ word-RAM
operations; a circuit must instead implement the dependence on address bits
with gates. For full APSP, the bound of Corollary~\ref{cor:apsp} is
weaker than the trivial $\Omega(m^2)$ needed to write $\Theta(m^2)$
outputs, so the informative statement there is the one about $m$ queries.
The exponent $\sigma$ is the unspecified constant from \cite{OAI}.

\paragraph{Status of the imported theorem.} Everything in this paper is
elementary given Theorem~\ref{thm:import}, which \cite{OAI} states without
hypotheses. We have not independently audited its proof. Result~06
\cite[Section~6]{R06} describes which parts of it have been checked.

\subsection{Outline}
In Section~\ref{sec:prelim} we state the imported theorem and show that
unambiguous verifiers give small tensor networks. Section~\ref{sec:lookup}
constructs the lookup graph, Section~\ref{sec:certificate} the unique
sorting certificate, and Section~\ref{sec:proof} combines them to prove
Theorem~\ref{thm:main}. Section~\ref{sec:decision} proves
Corollary~\ref{cor:variants}, Section~\ref{sec:apsp} proves
Corollary~\ref{cor:apsp}, and Section~\ref{sec:3sum} proves
Theorem~\ref{thm:3sum-intro}. We conclude with open problems in
Section~\ref{sec:open}.

\section{Preliminaries}\label{sec:prelim}

Let $[r]=\{0,\dots,r-1\}$. All logarithms are base two. For a predicate
$P$, $\ind{P}\in\{0,1\}$ is its indicator.

\paragraph{Tensor networks.} We follow \cite[Section~2.1]{OAI} and
\cite[Section~2]{R06}. Fix a field $k$. A network places a table
$\{0,1\}^u\to k$ at each vertex. Each wire pairs two argument positions,
and each unpaired position is a named exterior bit. The network represents
the function of its exterior bits obtained by multiplying vertex entries
and summing over all settings of the internal wires. It is
\emph{trivalent} if every table has arity at most three. Its
\emph{volume} counts vertices and all wires, including exterior wires.
Let $\vol_k(T)$ be the least volume of a trivalent network over $k$
representing $T$.

For $w\ge1$ let $q=2^w$ and $n_w=2qw$. Let $D_w:\{0,1\}^{n_w}\to\{0,1\}$
take $q$ words $X_0,\dots,X_{q-1}\in[q^2]$ of $2w$ bits each and return
$1$ exactly when they are pairwise distinct.

\begin{theorem}[{\cite[Theorem~2.2]{OAI}}, case $t=1$]\label{thm:import}
For every finite field $k$ there is $\sigma>0$ such that for all
sufficiently large $w$,
\[
\vol_k(D_w)>n_w\,w^{\sigma}=2q\,w^{1+\sigma}.
\]
\end{theorem}

The only other tool we need is the observation that a circuit with a
unique certificate yields a network for the function it verifies.

\begin{lemma}[Unambiguous verifiers give networks]\label{lem:verifier}
Let $V(x,c)$ be computed by a circuit with $s$ gates on $b$ input bits in
total, split into data bits $x$ and certificate bits $c$. Suppose that for
every $x$ the number of $c$ with $V(x,c)=1$ is $f(x)\in\{0,1\}$. Then over
every field, $f$ has a trivalent network with exterior bits $x$ and volume
at most $16s+5b+8$.
\end{lemma}
\begin{proof}
By \cite[Lemma~4.1]{R06}, which writes out the conversion in
\cite[proof of Proposition~2.3]{OAI}, there is a trivalent network of
volume at most $16s+4b+8$ representing $V$ with values $0,1$. It places the
relation table $\ind{c=g(a,b)}$ at each gate, ternary equality tables at
fan-outs, and the pin $\ind{c=1}$ at the output. For each input exactly one
internal assignment is consistent.

Close each certificate exterior bit with the unary table that is
identically $1$. This adds at most $b$ vertices and no wires, and sums
over $c$. The new network represents $\sum_cV(x,c)$, computed in the field.
That sum is the integer $f(x)\in\{0,1\}$, so it is $f(x)$ in every
characteristic.
\end{proof}

We emphasize that the sum over $c$ is a contraction. We do not claim that a
circuit can enumerate the certificates within this size;
Theorem~\ref{thm:import} bounds the volume of every network representing
$D_w$, however it is evaluated.

\section{Lookup by triangles}\label{sec:lookup}

We now construct the graph that turns triangles into table lookups. For
arrays $j=(j_i)_{i\in[q]}$, $x=(x_i)_{i\in[q]}$, and $y=(y_a)_{a\in[q]}$
of $w$-bit words, define
\[
R(j,x,y)=\bigwedge_{i\in[q]}\ind{x_i=y_{j_i}}.
\]
No promise is placed on the arrays; in particular $j$ need not be a
permutation.

\begin{lemma}[Lookup graph]\label{lem:graph}
For all $j,x,y$ there is a simple graph $G(j,x,y)$ on vertex set $[3q]$
with $3q$ edges, listed in a fixed order that is lexicographically sorted
for every input, such that:
\begin{enumerate}
\item for each $i$, a designated edge $e_i$ lies in a triangle if and only
if $x_i=y_{j_i}$, so $R(j,x,y)$ is the AND of the $\AEST$ outputs on
$e_0,\dots,e_{q-1}$;
\item every triangle contains exactly one vertex of $[q]$, and each vertex
of $[q]$ lies in at most one triangle;
\item $G(j,x,y)$ is 2-degenerate; and
\item the edge list is a projection of the bits of $j,x,y$ together with
constants.
\end{enumerate}
\end{lemma}
\begin{proof}
Name the vertices $I_i=i$, $A_a=q+a$, and $B_b=2q+b$ for $i,a,b\in[q]$,
and add the edges
\[
e_i=\{I_i,A_{j_i}\},\qquad f_i=\{I_i,B_{x_i}\},\qquad g_a=\{A_a,B_{y_a}\}.
\]
Within the $e$'s and within the $f$'s the endpoint $I_i$ distinguishes
the edges, and within the $g$'s the endpoint $A_a$ does. The three families
join different pairs of classes. So the graph is simple for every input.

The only edges are between classes $I$--$A$, $I$--$B$ and $A$--$B$, so
every triangle has one vertex in each class. Vertex $I_i$ has exactly the
two neighbors $A_{j_i}$ and $B_{x_i}$, so it lies in at most one triangle,
and
\[
e_i\text{ lies in a triangle}
\iff\{A_{j_i},B_{x_i}\}\text{ is an edge}
\iff x_i=y_{j_i}.
\]
This proves items~1 and~2.

List the edges as $e_0,f_0,e_1,f_1,\dots,e_{q-1},f_{q-1},g_0,\dots,g_{q-1}$.
Writing each edge with its smaller endpoint first, $e_i=(i,q+j_i)$ precedes
$f_i=(i,2q+x_i)$ because $q+j_i<2q$, the pairs for $i$ precede those for
$i+1$, and every $g_a=(q+a,2q+y_a)$ has smaller endpoint at least $q$. So
the list is sorted for every input.

For degeneracy, delete the $I$-vertices, each of degree $2$. Each
$A$-vertex then has degree $1$; delete them, and then the isolated
$B$-vertices.

Since $q$ is a power of two, the name $q+a$ is the bits of $a$ preceded by
fixed class bits, and likewise for $2q+b$. So the edge list is a projection.
\end{proof}

Hence an $\AEST$ circuit of size $S$ on these graphs yields a circuit for
$R$ of size $S+q-1$, plus $O(qw)$ constant gates for the fixed bits.

\section{A unique sorting certificate}\label{sec:certificate}

Next, we show how three lookups verify distinctness unambiguously. Write
each word of the distinctness input as $X_i=(a_i,b_i)$ with
$a_i,b_i\in[q]$, the first half most significant. A \emph{certificate}
consists of four arrays $j,z,u,v\in[q]^q$. Put $Y_c=(u_c,v_c)$ and let
$\id$ be the constant array $\id_i=i$. Define
\[
V(X;j,z,u,v)=R(j,\id,z)\wedge R(j,a,u)\wedge R(j,b,v)\wedge
\Inc(Y),
\qquad
\Inc(Y)=\bigwedge_{c=0}^{q-2}\ind{Y_c<Y_{c+1}},
\]
where $Y_c<Y_{c+1}$ compares $2w$-bit words. The intended certificate
takes $j_i$ to be the rank of $X_i$ in sorted order, $z=j^{-1}$, and $Y$
the sorted list.

\begin{lemma}[Uniqueness]\label{lem:unique}
For every $X$, the number of certificates accepted by $V(X;\cdot)$ is
$D_w(X)$.
\end{lemma}
\begin{proof}
Suppose $V$ accepts $(j,z,u,v)$. The first lookup says $z_{j_i}=i$ for
all $i$. If $j_i=j_{i'}$ then $i=z_{j_i}=z_{j_{i'}}=i'$, so $i\mapsto j_i$
is injective, hence a permutation of $[q]$, and $z$ is its inverse. The
other two lookups say $u_{j_i}=a_i$ and $v_{j_i}=b_i$, that is,
$Y_{j_i}=X_i$. Since $j$ is a permutation, $Y$ is a rearrangement of $X$,
and $\Inc$ says it is strictly increasing.

If $X$ has a repeated word, no rearrangement of $X$ is strictly increasing,
so no certificate is accepted. If $X$ is distinct, $Y$ must be the unique
increasing rearrangement of $X$; then $j_i$ must be the rank of $X_i$, and
$z$, $u$, $v$ are determined by $j$ and $Y$. This certificate passes every
check, so exactly one is accepted.
\end{proof}

Notice that neither a sorting circuit nor a permutation test is needed:
the inverse lookup forces a permutation, and the comparator chain checks
the proposed order. Both existence and uniqueness matter, since we will
sum over certificates in characteristic two.

\section{Proof of Theorem~\ref{thm:main}}\label{sec:proof}

We are now ready to prove our main result. It remains to pad the lookup
graph to every size $m$ and to account for the cost of the verifier.

\paragraph{Padding.} Let $m$ be large and let $q=2^w$ be the largest power
of two with $3q\le m$, so $m/6<q\le m/3$. For arrays $j,x,y$, let
$G_m(j,x,y)$ be $G(j,x,y)$ with a path of $p=m-3q$ edges attached at
$B_0=2q$, through the new vertices $3q,3q+1,\dots,m-1$:
\[
\{2q,3q\},\ \{3q,3q+1\},\ \dots,\ \{m-2,m-1\}.
\]
Append these edges to the list of Lemma~\ref{lem:graph}. The result has
exactly $m$ vertices and $m$ edges, and is simple. It is still sorted: the
first added edge has smaller endpoint $2q$, larger than that of every core
edge, and the later ones have increasing smaller endpoints. It is still
2-degenerate: delete the path from its far end inward, then proceed as
before. A pendant path lies on no triangle, so the outputs on
$e_0,\dots,e_{q-1}$ are unchanged. Since $L_m\ge w+2$, every name fits.

\paragraph{The verifier.} Let $C$ be a circuit with $S$ gates computing
$\AEST$ on $m$ edges, correct on the promised graphs. We build a verifier
for $D_w$: on input $X$ and certificate $(j,z,u,v)$, run $C$ on
$G_m(j,\id,z)$, $G_m(j,a,u)$ and $G_m(j,b,v)$, AND the outputs on the
designated edges, and AND the result with a comparator chain for
$\Inc(Y)$. The three edge lists are projections of the input bits plus
constants, and each comparison of $2w$-bit words takes $O(w)$ gates. The
verifier therefore has $s\le 3S+O(qw)$ gates and $b=2qw+4qw=6qw$ input
bits. By Lemma~\ref{lem:graph} it computes $V$, and by
Lemma~\ref{lem:unique} each $X$ has exactly $D_w(X)$ accepting
certificates.

\paragraph{Conclusion.} By Lemma~\ref{lem:verifier} over $\F_2$,
$D_w$ has a trivalent network of volume at most
$16(3S+O(qw))+30qw+8$. Theorem~\ref{thm:import} gives
\[
48S+O(qw)>2q\,w^{1+\sigma}.
\]
Since $w^\sigma\to\infty$, we get $S\ge q\,w^{1+\sigma}/48$ for large $w$.
As $q>m/6$ and $w=\log q\ge\log m-3$, this is
$\Omega(m(\log m)^{1+\sigma})$. Finally $\log B_m=\Theta(\log m)$, which
gives the bound in terms of input bits. \qed

\section{Counting, witnesses, and decision}\label{sec:decision}

In this section we prove Corollary~\ref{cor:variants}. The first two
parts follow from the structure of the lookup graph; the last two need one
more contraction.

\paragraph{Counting and per-edge witnesses.} By Lemma~\ref{lem:graph}(2),
$G_m(j,x,y)$ has exactly $\sum_i\ind{x_i=y_{j_i}}$ triangles, so
$R(j,x,y)=1$ exactly when the count equals $q$. A counting circuit
therefore replaces $C$ in the verifier, with $O(\log m)$ extra gates per
copy. A circuit that returns, for each edge, a triangle through it or a
marker determines every $\AEST$ output after $O(m\log m)$ gates of
decoding, which the bound absorbs. This proves parts~1 and~2.

\paragraph{Decision.} One cannot simply take the OR of the $\AEST$
outputs here, since that reduction goes the wrong way. Instead, we compute
$R$ from a decision circuit by summing over selector bits inside the
network.

Add a dummy vertex $B_*=3q$ with no neighbors in $A$. For selector bits
$s\in\{0,1\}^q$ let $G^s(j,x,y)$ replace each $f_i$ by
\[
f_i(s_i)=\begin{cases}\{I_i,B_{x_i}\},&s_i=1,\\ \{I_i,B_*\},&s_i=0,\end{cases}
\]
keeping the $e_i$ and $g_a$. This graph is simple, 2-degenerate by the same
order, and sorted in the same layout, since $f_i(s_i)$ still has smaller
endpoint $i$ and larger endpoint at least $2q>q+j_i$. Choosing each
endpoint of $f_i(s_i)$ costs $O(w)$ gates. Put $h_i=\ind{x_i=y_{j_i}}$.
Since $B_*$ has no $A$-neighbor, it lies on no triangle, and $G^s$ has a
triangle through $I_i$ exactly when $s_ih_i=1$. Its triangle-free
indicator is therefore $Q(s)=\prod_i(1-s_ih_i)$, and over every field
\[
\sum_{s\in\{0,1\}^q}(-1)^{|s|}\,Q(s)
=\prod_{i}\sum_{s_i\in\{0,1\}}(-1)^{s_i}(1-s_ih_i)
=\prod_i h_i=R(j,x,y).
\]
Over $\F_2$ the signs disappear. So if $T$ is a decision circuit with $S$
gates, the network of $\neg T$ on $G^s$, with its $q$ selector bits closed
by all-ones tables, represents $R$ with volume $O(S+qw)$. It uses one copy
of $T$, not $2^q$.

We use three such networks in place of the three lookups, sum over the
certificate as before, and include the comparator chain. Summing over the
selectors and over the certificate are both contractions, so the result
represents $\sum_{j,z,u,v}V=D_w$ with volume $O(S+qw)$. For general $m$,
attach the padding path at $B_0$ through the new vertices $3q+1,\dots,m$.
This gives $m$ edges on vertex set $[m+1]$ and preserves all the
properties above. Theorem~\ref{thm:import} then gives part~3. A circuit
that returns a triangle or a marker decides existence after $O(m\log m)$
gates of checking against the edge list, which gives part~4.

\section{Shortest paths}\label{sec:apsp}

We next prove Corollary~\ref{cor:apsp}. Here the reduction runs from
$\AEST$ to shortest paths, so Theorem~\ref{thm:main} transfers directly.

\begin{proof}[Proof of Corollary~\ref{cor:apsp}]
Given an $\AEST$ instance $G$ on $[m]$, form the layered graph $H$ with
vertices $u^0,u^1,u^2$ for $u\in[m]$ and, for each edge $\{u,v\}$ of $G$,
the four arcs $u^0\to v^1$, $v^0\to u^1$, $u^1\to v^2$, $v^1\to u^2$.
Name $u^\ell$ by two layer bits followed by the bits of $u$, so the arc
list is a projection of the edge list. Every path in $H$ goes from layer
$0$ to layer $1$ to layer $2$, so $H$ is acyclic and
$d_H(u^0,v^2)=2$ exactly when $u$ and $v$ have a common neighbor in $G$.
For an edge $\{u,v\}$ of $G$ this happens exactly when the edge lies in a
triangle. So the $m$ queries $(u^0,v^2)$, one per edge, give every $\AEST$
output after $O(\log m)$ gates each, and Theorem~\ref{thm:main} gives the
bound.
\end{proof}

The input has $\Theta(m\log m)$ bits, so this bound is again superlinear
in the input length. As noted in the introduction, for full APSP the
inherited bound is below the $\Omega(m^2)$ gates needed just to write the
output, and the informative statement is the one about $m$ queries. The
APSP hypothesis itself concerns dense weighted graphs, which we do not
address.

\section{3SUM}\label{sec:3sum}

Finally, we prove Theorem~\ref{thm:3sum-intro}. As discussed above, the
known reductions run from 3SUM to $\AEST$ \cite{VWX}, so we instead reduce
directly from distinctness.

\begin{definition}
For $n\ge3$ let $\ell_n=4\lceil\log n\rceil+2$. The function $\TSUM_n$
takes $n$ integers $z_1,\dots,z_n$ in $\ell_n$-bit two's complement, so
every $|z_i|\le n^4$ is representable, and returns $1$ exactly when
$z_a+z_b+z_c=0$ for some pairwise distinct positions $a,b,c$.
\end{definition}

The difficulty is that equal values must be detected only at
\emph{different} positions: placing $X_i$, $-X_i$ and $0$ in the instance
would create a solution for every $i$. We therefore encode the positions
into the integers.

\subsection{The reduction}
Let $M=4q$. Given $X_0,\dots,X_{q-1}\in[q^2]$, let $S(X)$ be the list of
the following $3q-1$ integers:
\begin{equation}\label{eq:3sum-red}
\begin{aligned}
a_i&=16(MX_i+i)+1, &&i\in[q],\\
b_k&=-16(MX_k+k)+2, &&k\in[q],\\
c_d&=16d-3, &&1\le d\le q-1.
\end{aligned}
\end{equation}

\begin{lemma}\label{lem:3sum-correct}
$\TSUM(S(X))=1-D_w(X)$.
\end{lemma}
\begin{proof}
\emph{A zero-sum triple has one integer of each type.} Modulo $16$ the
types are $a\equiv1$, $b\equiv2$, $c\equiv-3$. If a triple has $p$, $r$
and $t$ integers of the three types, then $p+r+t=3$ and
$p+2r-3t\equiv0\pmod{16}$. Since $-9\le p+2r-3t\le6$, in fact
$p+2r-3t=0$. Eliminating $p$ gives $r=4t-3$, so $t\ge1$, and $p=6-5t\ge0$
forces $t=1$. Hence $p=r=t=1$.

\emph{Such a triple encodes a collision.} By \eqref{eq:3sum-red},
\[
a_i+b_k+c_d=16\bigl(M(X_i-X_k)+(i-k+d)\bigr).
\]
Here $-(q-2)\le i-k+d\le 2q-2$, so $|i-k+d|<M$. The sum is zero exactly
when $X_i=X_k$ and $d=k-i$. As $1\le d\le q-1$, this requires $i<k$.
Conversely, if $X_i=X_k$ with $i<k$, then $d=k-i$ gives a zero-sum triple.
So $S(X)$ has a zero-sum triple exactly when $X$ has a repeated word.
\end{proof}

\begin{lemma}\label{lem:3sum-distinct}
The $3q-1$ integers in $S(X)$ are pairwise distinct for every $X$,
including lists with repeated words.
\end{lemma}
\begin{proof}
Different types differ modulo $16$. If $a_i=a_k$ then $M(X_i-X_k)=k-i$
with $|k-i|<M$, so $i=k$. The same holds for the $b$'s, and the $c$'s are
distinct by definition.
\end{proof}

So the reduction is valid even for the version of 3SUM whose input is
promised to be a set.

\begin{lemma}\label{lem:3sum-cost}
Every integer in $S(X)$ has absolute value less than $64q^3$, and $S$ is
computed by a circuit with $O(qw)$ gates.
\end{lemma}
\begin{proof}
From $X_i\le q^2-1$ and $i\le q-1$ we get $MX_i+i<4q^3$, which gives the
bound. As $M=2^{w+2}$ and $i<M$, the number $16(MX_i+i)$ is the
concatenation of the bits of $X_i$, of $i$ in $w+2$ bits, and of four
zeros. So each $a_i$ is a projection with constants. Each $b_k$ is a
two's complement negation plus a constant, which takes $O(w)$ gates.
The $c_d$ are constants.
\end{proof}

\subsection{The lower bound}
We restate Theorem~\ref{thm:3sum-intro} with the input format fixed.

\begin{theorem}\label{thm:3sum}
For every sufficiently large $n$, every circuit computing $\TSUM_n$ has
$\Omega(n(\log n)^{1+\sigma})$ gates. This holds even if the circuit is
required to be correct only on lists of pairwise distinct integers.
\end{theorem}
\begin{proof}
Let $w$ be largest with $3\cdot2^w-1\le n$, and put $q=2^w$. Then
$n<6q$. Put $R=64q^3$ and append the $n-3q+1$ padding integers
$3R+1,3R+2,\dots$ to $S(X)$. Every integer of $S(X)$ lies in $(-R,R)$. A
triple containing a padding integer has sum greater than $3R-2R>0$, so
padding creates no zero-sum triple. The padding integers exceed $R$ and
are distinct, so the padded list remains a set by
Lemma~\ref{lem:3sum-distinct}. Every entry has absolute value at most
$3R+n\le 8n^3\le n^4$ for large $n$, so it fits in $\ell_n$ bits.

Let $C$ be a circuit for $\TSUM_n$ with $s$ gates. By
Lemma~\ref{lem:3sum-correct}, $C$ preceded by this map and followed by one
NOT gate computes $D_w$. The map costs $O(qw)+O(n\ell_n)=O(n_w)$ gates.
By Lemma~\ref{lem:verifier} with an empty certificate and
Theorem~\ref{thm:import}, $16(s+O(n_w))+5n_w+8>n_ww^\sigma$, so
$s=\Omega(n_ww^\sigma)=\Omega(q(\log q)^{1+\sigma})$. Since $q>n/6$ and
$\log q\ge\log n-3$, this is $\Omega(n(\log n)^{1+\sigma})$.
\end{proof}

\begin{remark}
\begin{enumerate}
\item In bits, the input has $B=n\ell_n=\Theta(n\log n)$ bits and the
bound is $\Omega(B(\log B)^\sigma)$.
\item The three-list version, which asks for $a+b+c=0$ with $a\in A$,
$b\in B$, $c\in C$, inherits the bound directly from the lists
$\{a_i\},\{b_k\},\{c_d\}$; the type tags modulo $16$ are then
unnecessary. Negating the third list gives the version $a+b=c$.
\item The tags force the type of every zero-sum triple. Without them, a
triple such as $a_i+a_{i'}+c_d$ could vanish for reasons unrelated to
collisions.
\item As for $\AEST$, this is a circuit bound. It is far from the
conjectured $n^{2-o(1)}$ word-RAM time, and $\TSUM_n$ has circuits of size
$n^2\operatorname{polylog}n$.
\end{enumerate}
\end{remark}

\section{Open problems}\label{sec:open}
\begin{itemize}
\item \emph{Larger exponents.} Our bounds are superlinear only by a factor
of $(\log m)^{\sigma}$. Can tensor-network methods give an
$m^{1+\varepsilon}$ circuit lower bound for $\AEST$, for triangle
detection, or for 3SUM?
\item \emph{Dense inputs.} The hard instances are sparse. Is there a
superlinear-in-input bound, that is, $\Omega(v^2(\log v)^{\varepsilon})$
gates, for triangle detection from a $v\times v$ adjacency matrix, or for
APSP on dense weighted graphs?
\item \emph{Lookup as the hard core.} Our proof shows that three parallel
lookups plus a comparator chain suffice to verify distinctness
unambiguously. Which other problems embed parallel lookup at linear cost,
and so inherit the bound in the same way?
\end{itemize}

\paragraph{Acknowledgments.}
All bounds here are corollaries of \cite[Theorem~2.2]{OAI}, and all
credit for the underlying theorem belongs to its authors.

\begin{thebibliography}{9}
\bibitem{OAI}
OpenAI. The matrix multiplication exponent is nine fourths.
Preprint, November 24, 2026.
\href{https://github.com/openai/math/blob/3344da59991ed94f9ba711217b29c1d760badfa6/preprints/The-matrix-multiplication-exponent-is-nine-fourths-November-24-2026/paper.pdf}{Manuscript}.

\bibitem{R06}
C.~Cruz. Superlinear lower bounds for element distinctness and unambiguous
verification of collision from a tensor-network volume bound. Result~06 in
this collection, November 2026.

\bibitem{VWX}
V.~Vassilevska Williams and Y.~Xu. Monochromatic triangles, triangle
listing and APSP. In \emph{Proceedings of FOCS 2020}, pages 786--797, 2020.
\href{https://arxiv.org/abs/2007.09318}{arXiv:2007.09318}.
\end{thebibliography}
\end{document}
